% !TEX root = ../../hw03.tex
% Source: HW3 Intro to Math Analysis I (1).pdf, page 2, question 5.
% The source's quoted terms "sparse" and "incompleteness" are retained.
% Small covering length alone is not a criterion for metric incompleteness.
\begingroup
\renewcommand{\labelenumi}{(\alph{enumi})}
\begin{exercise}
  \label{ex:hw03:rational-interval-cover}
  \solutionlink{hw03:rational-interval-cover}
  Let the rational numbers in $[0,1]$ be enumerated by
  $q_0,q_1,\ldots,q_n,\ldots$. Let $\varepsilon>0$ and cover this set
  $Q=\{q_n\}$ by
  \[
    U_\varepsilon=\bigcup_{n=0}^{\infty}I_n(\varepsilon),
  \]
  \[
    I_n(\varepsilon)=(q_n-\varepsilon 2^{-n},q_n+\varepsilon 2^{-n}),
    \qquad n=0,1,2,\ldots.
  \]
  That is, $Q\subset U_\varepsilon$.
  \begin{enumerate}
    \item Show that the sum of the lengths of the intervals
      $\{I_n(\varepsilon)\}$ satisfies the bound
      \[
        \sum_{n=0}^{\infty}|I_n(\varepsilon)|\leq4\varepsilon.
      \]
    \item Since $\varepsilon>0$ can be arbitrarily small, the set $Q$ can
      be covered by intervals of arbitrarily small total length indicating
      that $Q$ is a rather ``sparse'' subset of $[0,1]$. In view of this
      fact, give your explanation regarding the ``incompleteness'' of $Q$
      in $[0,1]$.
  \end{enumerate}
\end{exercise}
\endgroup
